% Данный файл распространяетсяя под лицензией CC BY 4.0. Текст лицензии размещён на https://creativecommons.org/licenses/by/4.0/
% (c) Кафедра системного программирования, 2025

\documentclass[a4paper]{article}

\usepackage[a4paper, top=8mm, bottom=8mm, left=8mm, right=8mm]{geometry}

\usepackage{polyglossia}
\setdefaultlanguage[babelshorthands=true]{russian}
\setotherlanguage{english}

\usepackage{fontspec}
\setmainfont{FreeSerif}
\newfontfamily{\russianfonttt}[Scale=0.7]{DejaVuSansMono}

\usepackage[tiny, compact]{titlesec}

\usepackage{titling}
\setlength{\droptitle}{-1cm}
\pretitle{\begin{center}\begin{bfseries}\Large}
\posttitle{\par\end{bfseries}\end{center}}
\preauthor{\begin{center}\normalsize}
\postauthor{\par\end{center}\vspace{-1.8cm}}

\usepackage{hyperref}
\usepackage{bookmark}
\usepackage{csquotes}
% Сюда можно дописывать нужные вам пакеты.

\title{Изучение OAuth2: подключение провайдеров идентификации в приложение}

\author{Куляка Павел Дмитриевич}

% Дата много места занимает, её лучше не указывать.
\date{}

\begin{document}

\maketitle

\begin{flushright}
    Группа: \emph{25.Б41-мм}

    Кафедра: \emph{Системного программирования}

    % Степень/звание и должность научника мы лучше вас знаем, их можно не указывать.
    Научный руководитель: \emph{Смирнов Кирилл Константинович}

    Номер семестра практики: \emph{2}
\end{flushright} 


\section{Цель работы}

% Буквально парой предложений ввести в контекст, обязательно сформулировать цель работы. Если работа делается где-то глубоко внутри проекта, на введение можно потратить больше места, чтобы подвести читателя от сути проекта в целом к тому, что конкретно вам надо было сделать. Только без воды, никакого "в современном мире"!

Целью данной работы является изучение протоколов OAuth 2.0 и OpenID Connect, а также практическая реализация двух взаимодействующих приложений --- сервера авторизации (OIDC-провайдера) и клиента с возможностью входа через несколько внешних провайдеров.

% Это пояснение актуальности. Тут обязательно надо обрисовать полезность вашей работы для конечного пользователя (и заодно как-то описать, кто такой ваш конечный пользователь --- может, это соседняя команда). Даже если вы делаете работу "глубоко внутри", выясните у консультанта, чем это людям поможет. И опишите, кто выступал в роли "заказчика" вашей работы, т.е. кто и зачем предложил задачу. Можно без конкретных имён, но названия организаций вполне желательны.

Протоколы OAuth 2.0 и OIDC используются для авторизации пользователей через сторонние сервисы. Протокол OAuth 2.0 позволяет стороннему приложению получить ограниченный доступ к ресурсам пользователя без передачи ему пароля. OpenID Connect (OIDC) --- надстройка над OAuth 2.0, добавляющая слой аутентификации: сервер авторизации возвращает клиенту ID Token в формате JWT, содержащий информацию о пользователе. OIDC также определяет механизмы Discovery (автоматическое получение конфигурации сервера) и JWKS (публикация ключей для проверки подписи токенов). 

\section{Реализация}

\subsection{Архитектура}

% Раздел пишется в свободной форме, с минимумом технических подробностей (можно приводить ссылки на более подробные документы).

Проект состоит из двух независимых Flask-приложений: клиента и провайдера. Клиент выступает в роли OAuth-клиента и предоставляет пользователю интерфейс для входа через одного из поддерживаемых провайдеров. Провайдер реализует полноценный сервер авторизации OAuth 2.0 с расширением OpenID Connect.

% Вот это важно, опишите кратко, что использовали
% Реализация решения осуществлялась с использованием \dots (тут описать используемые технологии).

Реализация решения осуществлялась с использованием языка Python, веб-фреймворка Flask, библиотек PyJWT (работа с JWT-токенами), cryptography (криптографические операции), Werkzeug (хэширование паролей) и Requests (HTTP-запросы к провайдерам). Интерфейс свёрстан с помощью TailwindCSS.

% Тут надо кратко описать идею решения, спозиционировать его относительно предыдущих наработок\footnote{Если на них надо сослаться, то в подстраничной сноске, URL: \url{http://example.com} (дата обращения: не забываем указывать).} и существующих аналогов (тоже предельно кратко), по возможности обосновать принятые решения.

\subsection{Клиент}

Клиент поддерживает авторизацию через четырёх провайдеров:
\begin{itemize}
    \item Google --- протокол OpenID Connect (получение ID Token, проверка подписи через JWKS).
    \item GitHub --- протокол OAuth 2.0 (получение Access Token, запрос email через API).
    \item Yandex --- протокол OAuth 2.0 (аналогично GitHub).
    \item My Service --- собственный OpenID Connect-провайдер.
\end{itemize}

\subsubsection{Особенности подключения внешних провайдеров}

При подключении сторонних сервисов был учтён ряд особенностей. В экосистеме Яндекса существует два сервиса: Yandex ID (не поддерживает OIDC) и Yandex Identity Hub (поддерживает OIDC, PKCE, JWKS и Discovery). В работе был использован Yandex ID через OAuth 2.0 без OIDC. Также изначально вместо GitHub планировалось подключить авторизацию через Mail.ru, но, поскольку для регистрации приложения требовалось юридическое лицо, этот вариант был исключён.

\subsection{Провайдер}

Провайдер реализует полный процесс авторизации. Поддерживаются следующие endpoint'ы:
\begin{itemize}
    \item /oauth/.well-known/openid-configuration --- OpenID Connect Discovery.
    \item /oauth/authorize --- endpoint авторизации; создаёт код авторизации после успешной авторизации пользователя.
    \item /oauth/token --- token endpoint; обменивает код авторизации на Access Token и ID Token (JWT, подписанный RS256).
    \item /oauth/userinfo --- возвращает информацию о пользователе по действительному токену.
    \item /oauth/jwks.json --- публикует JWK Set для проверки подписи ID Token.
\end{itemize}

ID Token формируется как JWT. Подпись выполняется алгоритмом RS256 с использованием RSA-ключа длиной 2048 бит.

Для хранения данных (пользователи, зарегистрированные приложения, коды авторизации, токены) используется SQLite. Пароли пользователей хранятся в хэшированном виде (Werkzeug).

\section{Апробация}

\subsection{Автоматическое тестирование}

Для проверки корректности реализации написаны модульные тесты с использованием pytest. Покрыты следующие сценарии:

\begin{itemize}
    \item Авторизация через всех провайдеров (проверка формирования URL авторизации, корректность параметров state/nonce).
    \item Обмен кода авторизации на токены, в том числе с PKCE (plain и S256).
    \item Проверка подписи ID Token через JWKS.
    \item Валидация обязательных параметров и обработка ошибочных запросов (неверный client\_id, просроченный код, неверный \\redirect\_uri).
    \item Регистрация пользователей и OAuth-приложений на стороне провайдера.
    \item Управление приложениями (удаление, перегенерация client\_secret).
    \item Вспомогательные функции: распознавание email, извлечение email из структур данных, построение URL и заголовков.
\end{itemize}

\subsection{Интеграционная проверка}

Помимо автоматического тестирования была проведена ручная интеграционная проверка: клиент подключался к OIDC-провайдеру научного руководителя. Проверялся полный цикл --- редирект с клиента на провайдера, ввод логина и пароля на странице авторизации провайдера, получение авторизационного кода, редирект обратно на клиент, обмен кода на токены и валидация ID-токена (проверка подписи, issuer, audience и nonce). Всё отработало корректно, что подтверждает совместимость реализации со сторонними OIDC-серверами.

\subsection{Качество кода и CI}

Также настроен CI через GitHub Actions, включающий три workflow: запуск тестов (pytest), линтинг (ruff) и проверка типов (mypy).

\section{Выездной день}
В рамках учебной практики состоялся выездной день в офис компании YADRO. Основной целью мероприятия было знакомство студентов с деятельностью компании, её продуктами и возможностями профессионального развития для молодых специалистов.

Представители компании рассказали о ключевых направлениях работы YADRO, о разрабатываемых продуктах и о задачах, которые решают инженерные команды. Особенно интересной для меня оказалась тема систем хранения данных: было полезно узнать, как обеспечиваются надёжность, производительность и отказоустойчивость таких систем.

Лекции были построены не только как обзор деятельности компании, но и как небольшое введение в реальные инженерные задачи. Спикеры объясняли основные принципы работы технологий и делились практическими примерами. Отдельное внимание уделялось вопросам измерения производительности: рассказывали о выборе метрик, поиске причин снижения эффективности и способах устранения «узких мест» в системе.

Также нам подробно рассказали о возможностях для студентов. Особое внимание было уделено стажировке «Импульс», на которую можно попасть уже со второго курса, а также другие направления для начинающих специалистов. Кроме того, сотрудники упомянули практические курсы и образовательные программы, которые компания проводит для студентов в течение года.

В конце мероприятия прошла встреча со стажёрами компании, среди которых были выпускники и студенты математико-механического факультета СПбГУ. Они поделились своим опытом работы и обучения, рассказали о проектах, в которых участвуют, и ответили на вопросы о стажировках и работе в компании. Было интересно услышать, как знания, полученные в университете, применяются на практике.

Выездной день оставил положительное впечатление. Мероприятие помогло лучше представить, как устроена работа в крупной IT-компании, и дало более понятное представление о возможных направлениях дальнейшего профессионального развития.

\section{Заключение}

В ходе выполнения работы получены следующие результаты:

\begin{itemize}
    \item Реализован сервер авторизации (провайдер) OAuth 2.0 / OpenID Connect;
    \item Реализовано клиентское веб-приложение с поддержкой входа через Google (OIDC), GitHub (OAuth 2.0), Яндекс (OAuth 2.0) и локальный OIDC-провайдер;
    \item Разработана система управления пользователями и приложениями (регистрация, авторизация, удаление приложений, сброс секретов);
    \item Написаны модульные тесты , настроена непрерывная интеграция через GitHub Actions;
    \item Проведена интеграционная проверка с OIDC-провайдером научного руководителя, подтвердившая совместимость реализации;
    \item Реализована поддержка множества внешних провайдеров через единый интерфейс конфигурации, что позволяет легко добавлять новых провайдеров путём описания их параметров в конфигурационном словаре.
\end{itemize}

Репозиторий: \url{https://github.com/pavelkuliaka/educational-practice}

\end{document}
